\documentclass[10pt,a4paper]{article}
\usepackage{amsmath,amssymb,amsthm,graphicx}
% Card-local minimal preamble (journal unit sigma-2025-018-mass-from-an-extrinsic-point-of-view, P2-ahl16 precedent).
%% Verbatim from the pinned source: the author theorem environments and macro block,
%% so every macro and environment of the counted spans is available.
% Omitted and why: the documentclass line (the journal class sigma is not on the host;
%% the assembler supplies the class) and page-geometry/font blocks (they fight the
%% forced xelatex engine).
\usepackage{amsmath,amssymb,amsthm,mathtools}
\usepackage{xcolor}
\usepackage{url}
\usepackage[hidelinks]{hyperref}
% Patch 33: the source no-break space (U+00A0) is treated as an ordinary space.
\catcode"00A0=10 \relax
\usepackage{xfrac}

\numberwithin{equation}{section}
\newtheorem{Theorem}{Theorem}[section]
\newtheorem*{Theorem*}{Theorem}
\newtheorem{Corollary}[Theorem]{Corollary}
\newtheorem{Lemma}[Theorem]{Lemma}
\newtheorem{Proposition}[Theorem]{Proposition}
\newtheorem{Conjecture}[Theorem]{Conjecture}
 { \theoremstyle{definition}
\newtheorem{Definition}[Theorem]{Definition}
\newtheorem{Note}[Theorem]{Note}
\newtheorem{Example}[Theorem]{Example}
\newtheorem{Remark}[Theorem]{Remark}
\newtheorem{Question}[Theorem]{Question}}
\usepackage{cleveref}
\crefname{originalreference}{original reference}{original references}
\Crefname{originalreference}{Original reference}{Original references}
\makeatletter
\newlabel{decay gradient}{{2.8}{source}{}{sourcepointer-1}{}}
\newlabel{decay gradient@cref}{{[originalreference][0][]2.8}{[1][0][]source}{}{}{}}
\makeatother
\crefname{originalreference}{original reference}{original references}
\Crefname{originalreference}{Original reference}{Original references}
\makeatletter
\newlabel{div Y interno G}{{3.4}{source}{}{sourcepointer-2}{}}
\newlabel{div Y interno G@cref}{{[originalreference][0][]3.4}{[1][0][]source}{}{}{}}
\makeatother
\crefname{originalreference}{original reference}{original references}
\Crefname{originalreference}{Original reference}{Original references}
\makeatletter
\newlabel{equation relation G and T}{{3.2}{source}{}{sourcepointer-3}{}}
\newlabel{equation relation G and T@cref}{{[originalreference][0][]3.2}{[1][0][]source}{}{}{}}
\makeatother
\crefname{originalreference}{original reference}{original references}
\Crefname{originalreference}{Original reference}{Original references}
\makeatletter
\newlabel{gradient field}{{2.8}{source}{}{sourcepointer-4}{}}
\newlabel{gradient field@cref}{{[originalreference][0][]2.8}{[1][0][]source}{}{}{}}
\makeatother
\crefname{originalreference}{original reference}{original references}
\Crefname{originalreference}{Original reference}{Original references}
\makeatletter
\newlabel{trace Lovelock}{{2.6}{source}{}{sourcepointer-5}{}}
\newlabel{trace Lovelock@cref}{{[originalreference][0][]2.6}{[1][0][]source}{}{}{}}
\makeatother
\crefname{originalreference}{original reference}{original references}
\Crefname{originalreference}{Original reference}{Original references}
\makeatletter
\newlabel{trace T}{{3.1}{source}{}{sourcepointer-6}{}}
\newlabel{trace T@cref}{{[originalreference][0][]3.1}{[1][0][]source}{}{}{}}
\makeatother
\title{The mass of an asymptotically flat manifold from an extrinsic point of view}
\author{Girao, Frederico, de Sousa, Alexandre}
\date{doi:10.3842/SIGMA.2025.018; CC BY 4.0}
\begin{document}
\maketitle
\setcounter{section}{3}
\setcounter{Theorem}{5}
\noindent\textit{Editorial note.} All delivered mathematical text is the author's own, copied from the pinned journal source file sources/sigma-2025-018-mass-from-an-extrinsic-point-of-view/pinned-original-source.tex (Mass from an Extrinsic Point of View); no statement or proof step was rewritten, reordered or removed.\par\medskip
\noindent\textit{Editorial note.} The counted proof is the author's complete argument: All delivered mathematical text is the author own, copied verbatim from the pinned journal source. Statement source lines 480-498 with the author own proof at 500-524 (label `main theorem`). Printed number 3.1 from the source walk; verified against the host aux. Pre-dispatch self-certification: DOI-wide key scan found no other source carrying this face.\par\medskip
\noindent\textit{Editorial note.} The label main theorem names the counted statement, and the proof environment immediately following it in the source is delivered verbatim as its proof.\par\medskip
\noindent\textit{Editorial note.} The counted statement is carried by the author's own Theorem environment in the article source.\par\medskip
\noindent\textit{Editorial note.} Journal-unit assembly: card-local preamble (author theorem environments and macros verbatim), article class, bibliography backend none; any literature citation carried by the delivered span is disclosed in the notes of this card.\par\medskip
\noindent\textit{Editorial note.} Author names were transliterated to ASCII for the card face by stripping diacritics (1 name(s) affected); the SOURCE.json metadata keeps the original accented form.\par\medskip
\noindent\textit{Editorial scope.} the counted statement and its own proof environment only\par\medskip
\noindent\textit{Allowed standard prerequisites (editorial).}\begin{itemize}
\item The ADM mass and the Gauss-Bonnet-Chern mass of an asymptotically f.... Supplies background that the delivered proof uses. The entry is declared as a dependency; it is not re-proved on the card, and the load-bearing steps of the counted statement are carried out inside the delivered span.
\end{itemize}\medskip
\ifdefined\hypertarget\hypertarget{sourcepointer-1}{}\fi
\noindent\textit{Source reference.} Original reference 2.8: sources/sigma-2025-018-mass-from-an-extrinsic-point-of-view/pinned-original-source.tex, source lines 289--289. Navigation pointer invoked at line(s) 505: the counted text refers to this environment of the same journal unit, which is not retained in this card, so the reference is declared as a pointer carrying the number the article itself prints for it (the section and the place of the statement inside that section, derived from the shared section-scoped theorem counter of the source); the cited item is an allowed input and none of it is reproduced inside the span.\par\medskip
\ifdefined\hypertarget\hypertarget{sourcepointer-2}{}\fi
\noindent\textit{Source reference.} Original reference 3.4: sources/sigma-2025-018-mass-from-an-extrinsic-point-of-view/pinned-original-source.tex, source lines 415--415. Navigation pointer invoked at line(s) 517: the counted text refers to this environment of the same journal unit, which is not retained in this card, so the reference is declared as a pointer carrying the number the article itself prints for it (the section and the place of the statement inside that section, derived from the shared section-scoped theorem counter of the source); the cited item is an allowed input and none of it is reproduced inside the span.\par\medskip
\ifdefined\hypertarget\hypertarget{sourcepointer-3}{}\fi
\noindent\textit{Source reference.} Original reference 3.2: sources/sigma-2025-018-mass-from-an-extrinsic-point-of-view/pinned-original-source.tex, source lines 381--381. Navigation pointer invoked at line(s) 523: the counted text refers to this environment of the same journal unit, which is not retained in this card, so the reference is declared as a pointer carrying the number the article itself prints for it (the section and the place of the statement inside that section, derived from the shared section-scoped theorem counter of the source); the cited item is an allowed input and none of it is reproduced inside the span.\par\medskip
\ifdefined\hypertarget\hypertarget{sourcepointer-4}{}\fi
\noindent\textit{Source reference.} Original reference 2.8: sources/sigma-2025-018-mass-from-an-extrinsic-point-of-view/pinned-original-source.tex, source lines 289--289. Navigation pointer invoked at line(s) 505: the counted text refers to this environment of the same journal unit, which is not retained in this card, so the reference is declared as a pointer carrying the number the article itself prints for it (the section and the place of the statement inside that section, derived from the shared section-scoped theorem counter of the source); the cited item is an allowed input and none of it is reproduced inside the span.\par\medskip
\ifdefined\hypertarget\hypertarget{sourcepointer-5}{}\fi
\noindent\textit{Source reference.} Original reference 2.6: sources/sigma-2025-018-mass-from-an-extrinsic-point-of-view/pinned-original-source.tex, source lines 203--203. Navigation pointer invoked at line(s) 523: the counted text refers to this environment of the same journal unit, which is not retained in this card, so the reference is declared as a pointer carrying the number the article itself prints for it (the section and the place of the statement inside that section, derived from the shared section-scoped theorem counter of the source); the cited item is an allowed input and none of it is reproduced inside the span.\par\medskip
\ifdefined\hypertarget\hypertarget{sourcepointer-6}{}\fi
\noindent\textit{Source reference.} Original reference 3.1: sources/sigma-2025-018-mass-from-an-extrinsic-point-of-view/pinned-original-source.tex, source lines 366--366. Navigation pointer invoked at line(s) 523: the counted text refers to this environment of the same journal unit, which is not retained in this card, so the reference is declared as a pointer carrying the number the article itself prints for it (the section and the place of the statement inside that section, derived from the shared section-scoped theorem counter of the source); the cited item is an allowed input and none of it is reproduced inside the span.\par\medskip
\begin{Theorem} \label{main theorem}
 Let $n$ and $q$ be integers such that $n \geq 3$ and $0 < q < n/2$.
 If $\psi\colon  M^n \hookrightarrow \bigl(\mathbb{R}^d, \bar{\delta}\bigr)$, $d > n$, is an
 asymptotically Euclidean immersion of order $\tau > \tau_q$ such that
 \(L_{(q)}\), the $q$-th Gauss-Bonnet curvature of the induced metric
 \(\psi^{\ast}\bar{\delta}\), is integrable, then the functions $S_{2q}$
 and $\bigl\langle S_{2q+1},\bar{Z} \bigr\rangle$ are integrable
 and the $q$-th Gauss--Bonnet--Chern mass of $\bigl(M, \psi^\ast\bar{\delta}\bigr)$ satisfies
 the following integral identity:
 \begin{equation*}
  \mathrm{m}_q\bigl(M,\psi^*\bar{\delta}\bigr) = a(n,q) \biggl[
   (n-2q) \int_M S_{2q} \, {\rm d}M
   + (2q+1) \int_M \langle S_{2q+1}, \bar{Z} \rangle \, {\rm d}M \biggr],
 \end{equation*}
 where \({\rm d}M\) is the Riemannian measure on \(\bigl(M, \psi^{\ast}\bar{\delta}\bigr)\) and
 \begin{equation*}
  a(n,q) = \frac{(2n)!(n-2q-1)!}{2^q(n-1)! \, \omega_{n-1}}.
 \end{equation*}
\end{Theorem}

\begin{proof}
 Let $E$ be one of the ends of $M$. Consider the function $f\colon M \to \mathbb{R}$, with
 \begin{equation*}
  f = \frac{\psi^\ast \bar{\rho}^2}{2}.
 \end{equation*}
 Since $\psi\colon  M^n \looparrowright \mathbb{R}^d$ is an asymptotically Euclidean immersion of order $\tau > \tau_q$, the function $f$ satisfies (\ref{decay gradient}), and hence, by Corollary \ref{gradient field},
 \begin{equation} \label{eq 1}
  \mathrm{m}_q \bigl(E,\psi^\ast \bar{\delta}\bigr) = -b(n,q) \lim_{\rho \to \infty} \int_{S_\rho} G_{(q)}(\nabla f,\nu)\, {\rm d}S_\rho,
 \end{equation}
 where the geometric quantities in the integral are computed with respect
 the induced metric~\(\psi^\ast \bar{\delta}\).

 Note that
 \begin{equation*}
  \nabla f = \psi^\ast \bar{Z}^\top.
 \end{equation*}
 Thus, applying equation (\ref{eq 1}) to all the ends of $M$ together
 with the divergence theorem and identity (\ref{div Y interno G}), we find
 \begin{equation} \label{eq 2}
  \mathrm{m}_q\bigl(M,\psi^\ast \bar{\delta}\bigr)
  = a(n,q) \int_M \bigl[ (n-2q)S_{2q} + (2q + 1) \bigl\langle S_{2q+1}, \bar{Z} \bigr\rangle \bigr] {\rm d}M.
 \end{equation}

 It remains to show that both $S_{2q}$ and $\bigl\langle S_{2q+1}, \bar{Z} \bigr\rangle$ are integrable, so the right hand side of~(\ref{eq 2}) can be broken in two. The integrability of $S_{2q}$ follows from (\ref{trace Lovelock}), (\ref{trace T}), (\ref{equation relation G and T}) and the integrability of $L_{(q)}$. Once we know that $S_{(2q)}$ is integrable, the integrability of $\bigl\langle S_{2q+1}, \bar{Z} \bigr\rangle$ follows from the fact that the left hand side of (\ref{eq 2}) is finite.
\end{proof}

\end{document}
